\documentclass{report}
\usepackage{amsmath,amssymb}
\setlength{\parindent}{0mm}
\setlength{\parskip}{1em}
\begin{document}
\begin{center}
\rule{6in}{1pt} \
{\large Srikanth Allu \\
{\bf Development of single domain framework to model 3D Li-Ion intercalation batteries using the Advanced Multi Physics (AMP) package. }}

OAK RIDGE NATIONAL LABORATORY \\ PO BOX 2008 MS6164 \\ OAK RIDGE TN 37831-6164
\\
{\tt allus@ornl.gov}\\
Sreekanth Pannala \\
Partha Mukherjee\\
Bobby Philip\\
John Turner\end{center}

Modeling of batteries involves simulation of coupled charge and thermal
transport, interfacial electrochemical reactions across the porous 3D
structure of the electrodes (cathodes and anodes) and electrolyte system.
The predominant approach has been the pseudo-2D modeling based on the
porous electrode theory where transport in the solid-phase is considered
in spherical active particles and the diffusion of Li is solved
separately at each point along the thickness direction. These models are
limited in their ability to capture spatial variations or able to handle
multidimensional structure of the electrodes and do not scale to the
system-level. In our approach we adapt the equations based on rigorous
homogenization procedures which would allow modeling 3D electrodes. The
resulting equations become stiff differential algebraic systems on a
single domain and require a robust solution methodology. These
formulations are implemented in AMPERES code, an external package for AMP
software. In the talk we demonstrate the importance of the coupled
solution procedure using Newton methods and accelerate using Krylov
iterative solvers. Experiments based on various preconditioning
strategies are presented.


\end{document}
